\documentclass[11pt]{article}
\usepackage[margin=1in]{geometry}
\usepackage{amsmath,amssymb}
\usepackage{booktabs}
\usepackage{tabularx}
\usepackage{enumitem}
\usepackage{xcolor}
\usepackage[most]{tcolorbox}
\usepackage[hidelinks]{hyperref}

\newtcolorbox{analogy}{colback=blue!4,colframe=blue!45!black,
  fonttitle=\bfseries,title=Hardware analogy,breakable}
\newtcolorbox{pitfall}{colback=red!4,colframe=red!50!black,
  fonttitle=\bfseries,title=Common misconception,breakable}

\title{README: A Primer for the MSARCH Project\\[4pt]
\large Bioinformatics, Co-Expression Networks, QUBO, QAOA, and cuQuantum\\
for readers who know SIMD, assembly, and CUDA}
\date{September 2026}

\begin{document}
\maketitle

\section*{How to use this document}

The main specification (\texttt{main.tex}) is written for readers who already
know transcriptomics and WGCNA. This primer assumes only that you know
computer architecture: SIMD/SIMT execution, memory hierarchies, assembly-level
reasoning, and CUDA. Every biological or quantum concept is reduced to the
arrays and operations you would actually write kernels for. Read it once
front to back; afterwards, Section~\ref{sec:map} maps each pipeline module to
the concepts here, and Section~\ref{sec:glossary} is a glossary.

\tableofcontents

%======================================================================
\section{Biology in one page}

\subsection{Genes, RNA, and expression}

Every cell in your body carries the same DNA, a string over the alphabet
$\{A,C,G,T\}$ of about $3\times10^9$ characters. A \textbf{gene} is a
region of that string that encodes a product, usually a protein. Humans have
about 20{,}000 protein-coding genes.

Cells do not use all genes at once. To use a gene, the cell
\textbf{transcribes} it into a temporary copy made of RNA (messenger RNA,
mRNA), and ribosomes \textbf{translate} that copy into protein. This is the
``central dogma'': DNA $\to$ RNA $\to$ protein.

\begin{analogy}
DNA is the program image on disk: every cell has the full binary. RNA
transcripts are the instructions currently paged into the instruction cache.
Counting how many copies of each gene's RNA are present tells you which
``functions'' the cell is calling right now, and how hot each one is. The
count for a gene is its \textbf{expression level}.
\end{analogy}

\subsection{Transcriptomics and RNA-Seq}

The \textbf{transcriptome} is the full set of RNA in a tissue sample;
\textbf{transcriptomics} is measuring it. The dominant technology,
\textbf{RNA-Seq}, shreds the RNA, sequences millions of short fragments
(``reads''), aligns each read back to the genome, and counts how many reads
land in each gene.

The result is a plain integer matrix:
\[
C \in \mathbb{N}^{G_0\times S}, \qquad C_{is} = \text{reads from gene } i
\text{ in sample } s.
\]
Rows are genes (tens of thousands), columns are samples (tissue from different
people). In this project, each column is a piece of brain from one deceased
donor. That matrix is the only input the GPU ever sees; everything biological
is metadata attached to rows (gene names) and columns (donor diagnosis, age,
etc.).

\begin{pitfall}
``Bulk'' RNA-Seq, used here, averages over all cells in a tissue chunk
(neurons, immune cells, support cells). Single-cell RNA-Seq gives one column
per cell instead, and is out of scope. A change in a bulk expression value can
mean a gene changed \emph{or} that the mix of cell types changed (e.g., neurons
died). Biologists interpret results with this in mind.
\end{pitfall}

\subsection{Why raw counts need preprocessing}

Raw counts are not directly comparable:
\begin{itemize}[nosep]
  \item \textbf{Library size:} one sample may simply have been sequenced
        twice as deeply, doubling every count. Normalization (CPM, TMM)
        rescales each column.
  \item \textbf{Skew:} counts span orders of magnitude, so we take
        $\log_2(\text{count}+1)$ to compress the range.
  \item \textbf{Batch effects:} samples processed on different days or
        machines differ systematically. These are removed by regressing out
        known technical variables (batch, RNA quality ``RIN'', time between
        death and tissue freezing ``PMI'').
  \item \textbf{Filtering:} genes that are nearly always zero carry noise, not
        signal, and are dropped. We then keep the $G$ most variable genes.
\end{itemize}
This is Module~1. It is cheap ($O(GS)$) and mostly done on the CPU.

\subsection{Alzheimer's disease in brief}

Alzheimer's disease (AD) is the most common cause of dementia. Its brain
hallmarks are \textbf{amyloid plaques} (clumps of a protein fragment between
cells) and \textbf{tau tangles} (twisted protein fibers inside neurons).
Pathologists score these on standard scales: \textbf{Braak stage} (how far
tangles have spread, 0--VI) and \textbf{CERAD} (plaque density). Late-onset AD
is \textbf{polygenic}: many genes each nudge risk slightly (the largest single
factor is the \textit{APOE}~$\varepsilon4$ variant). That is why we look at
genes as a network instead of one at a time.

\subsection{The ROSMAP dataset}

ROSMAP (Religious Orders Study and Rush Memory and Aging Project) follows
older adults who agree to yearly cognitive testing and brain donation at
death. For each donor we get brain RNA-Seq plus clinical diagnosis, cognitive
test trajectories, and pathology scores. The data live on the AD Knowledge
Portal (Synapse) and require registering and signing a data use certificate
before download; plan for that delay. Public data from NCBI GEO can be used
meanwhile.

%======================================================================
\section{From expression to co-expression networks}

\subsection{The core idea}

If two genes rise and fall together across hundreds of people, they are
likely controlled by the same mechanism or work in the same process. Groups of
such genes are called \textbf{modules}. If a module's overall activity tracks
disease severity, the module (and its most central genes) are candidates for
understanding or treating the disease. That is all of WGCNA in one paragraph;
the rest is linear algebra.

\subsection{Correlation is a dot product}

Take gene $i$'s row $x_i\in\mathbb{R}^S$. Subtract its mean and divide by its
norm to get a unit vector $z_i$. Then the Pearson correlation between genes
$i$ and $j$ is just
\[
r_{ij} = z_i\cdot z_j = \sum_{s=1}^{S} z_{is} z_{js} \in [-1,1].
\]
Stacking all standardized rows into $Z\in\mathbb{R}^{G\times S}$, every pair at
once is
\[
R = Z Z^{\top}.
\]

\begin{analogy}
This is a GEMM with inner dimension $S$ (hundreds) and output $G\times G$
(up to $20{,}000^2 = 4\times10^8$ entries). It is exactly the tiled
shared-memory matrix multiply from any CUDA course, and because $R$ is
symmetric you only need the upper triangle (BLAS \texttt{SYRK}). Arithmetic
intensity is high, so this stage is compute-bound.
\end{analogy}

\subsection{Soft thresholding: turning correlation into a network}

A network needs edge weights $a_{ij}\ge0$. WGCNA uses
\[
a_{ij} = |r_{ij}|^{\beta}, \qquad \beta \approx 6\text{--}14.
\]
Raising to a power keeps strong correlations (0.9$^{6}\approx0.53$) and
crushes weak ones (0.3$^{6}\approx0.0007$) without a hard cut-off. $\beta$ is
chosen so that the network's degree distribution looks \textbf{scale-free}
(a few highly connected hubs, many sparsely connected nodes), which is
empirically what biological networks look like.

\begin{analogy}
A pure element-wise map over $G^2$ values: a textbook memory-bound kernel.
Every element is read once and written once; the only optimizations that
matter are coalescing, vector loads (\texttt{float4}), and fusing it with the
next pass (row sums).
\end{analogy}

The \textbf{connectivity} (weighted degree) of gene $i$ is the row sum
$k_i = \sum_u a_{iu}$.

\subsection{Topological Overlap: counting shared neighbors}

Correlation is noisy with a few hundred samples. WGCNA makes the similarity
more robust by asking not only ``are $i$ and $j$ connected?'' but ``do $i$
and $j$ connect to the same other genes?''
\[
l_{ij} = \sum_{u} a_{iu} a_{uj} \quad\Longleftrightarrow\quad L = A\cdot A,
\]
\[
\text{TOM}_{ij} = \frac{l_{ij} + a_{ij}}{\min(k_i,k_j) + 1 - a_{ij}}.
\]
$l_{ij}$ is the weighted count of 2-hop paths from $i$ to $j$. The
denominator normalizes it into $[0,1]$. The distance used for clustering is
$D = 1 - \text{TOM}$.

\begin{analogy}
$A\cdot A$ with $G=20{,}000$ is $2G^3 = 1.6\times10^{13}$ FLOPs as a plain
GEMM, or half that with SYRK because $A$ is symmetric. This one matrix
product is the reason the project exists: it is the $O(G^3)$ term. The TOM
formula afterward is another memory-bound element-wise pass that needs
$k_i$ and $k_j$, which you can cache in shared memory per tile.
\end{analogy}

\begin{pitfall}
$O(G^3)$ is \emph{cubic}, not exponential. Doubling $G$ costs $8\times$ in
time and $4\times$ in memory. The real-world pain is memory: a
$20{,}000\times20{,}000$ float matrix is 1.6~GB, and you need two or three of
them. CPU tools often split genes into blocks to fit, which throws away
between-block overlap. A GPU with 8+~GB holds the whole thing.
\end{pitfall}

\subsection{Clustering into modules}

With the distance matrix $D$, standard \textbf{hierarchical clustering}
(average linkage) builds a binary tree (dendrogram) by repeatedly merging the
two closest clusters. \textbf{Dynamic Tree Cut} then cuts the tree into
branches of varying height; each branch with at least $\sim$30 genes becomes a
module. Genes that fit nowhere go to a ``grey'' leftover module (label 0).
This step is irregular, pointer-chasing, and sequential in nature, which is why
it stays on the CPU in our baseline---and why it is the natural place to plug
in alternative algorithms.

\subsection{Eigengenes, hubs, and linking to disease}

\begin{itemize}[nosep]
  \item \textbf{Module eigengene:} the first principal component of the
        module's rows of $Z$ (the top right singular vector from an SVD).
        Think of it as the module's ``average'' expression pattern across
        people---one number per donor per module.
  \item \textbf{Module--trait correlation:} correlate each eigengene with
        disease variables (diagnosis, Braak, CERAD, cognitive decline). A
        module whose eigengene tracks disease is ``AD-associated.''
  \item \textbf{Hub genes:} genes in that module most correlated with its
        eigengene ($k_{\text{ME}}$). These are the candidate biomarkers.
  \item \textbf{Enrichment:} check whether a module is over-represented for
        known functions (e.g., immune response) or cell types (e.g.,
        microglia, the brain's immune cells) using a hypergeometric test.
        This is how a list of gene names becomes a biological story.
\end{itemize}

%======================================================================
\section{Optimization problems: QUBO and Ising}

\subsection{QUBO}

A \textbf{Quadratic Unconstrained Binary Optimization} problem is
\[
\min_{x\in\{0,1\}^n} \; x^{\top} Q x = \sum_{i} Q_{ii}x_i + \sum_{i<j}
(Q_{ij}+Q_{ji}) x_i x_j,
\]
for a real $n\times n$ matrix $Q$ (linear terms sit on the diagonal because
$x_i^2 = x_i$). Many NP-hard problems (max-cut, graph partitioning, scheduling)
can be encoded this way; constraints are added as penalty terms.

\begin{analogy}
Brute force is a loop over all $2^n$ bit patterns of an $n$-bit register,
evaluating a quadratic form each time: embarrassingly parallel on a GPU, one
thread per pattern, and practical up to $n\approx 30$--$35$. We use exactly
this as the ground-truth baseline for the quantum solver.
\end{analogy}

\subsection{Ising form}

Substituting $x_i = (1+s_i)/2$ with spins $s_i\in\{-1,+1\}$ turns a QUBO into
an Ising energy $\sum_i h_i s_i + \sum_{i<j} J_{ij} s_i s_j$ (plus a constant).
Quantum algorithms work natively in this form because a qubit measured in the
computational basis yields $\pm1$ under the Pauli-$Z$ operator.

\subsection{Module detection as a QUBO}

To split a weighted graph into two groups, assign each node a spin. Newman's
\textbf{modularity} measures how much more edge weight falls inside groups
than expected at random:
\[
Q = \frac{1}{4m}\, s^{\top} B s, \qquad B = W - \frac{k k^{\top}}{2m},
\]
where $W$ is the weight matrix, $k$ the degrees, and $2m=\sum k_i$.
Maximizing $Q$ is minimizing $-s^{\top}Bs$: an Ising problem with
$J=-B$. Recursive splitting, or a one-hot encoding with one bit per
(node, community) pair, gives more than two groups.

\section{Quantum computing, just enough}

\subsection{Qubits and the state vector}

A classical $n$-bit register holds one of $2^n$ values. An $n$-qubit
register's state is a vector of $2^n$ complex \textbf{amplitudes}
$\psi\in\mathbb{C}^{2^n}$ with $\|\psi\|_2=1$. Measuring it returns bit
pattern $b$ with probability $|\psi_b|^2$ and destroys the superposition.

\begin{analogy}
Simulating $n$ qubits on classical hardware means storing $\psi$ explicitly:
$2^n$ complex numbers. At complex64 (8 bytes) that is 8~GiB for 30 qubits and
16~GiB for 31. Every added qubit doubles memory. This is the hard wall for
the quantum part of the project and the reason it runs on a small, reduced
graph (a few dozen module eigengenes), never on 20{,}000 genes.
\end{analogy}

\subsection{Gates are structured matrix--vector products}

A quantum program is a sequence of \textbf{gates}, each a unitary matrix
applied to $\psi$. A 1-qubit gate is a $2\times2$ matrix applied to qubit $q$;
on the full state it acts on all $2^{n-1}$ pairs of amplitudes whose indices
differ only in bit $q$.

\begin{analogy}
Applying a 1-qubit gate is a strided butterfly over the state vector, like
one stage of an FFT: pair up elements $i$ and $i \oplus (1\ll q)$ and apply a
$2\times2$ complex multiply to each pair. It is memory-bandwidth-bound and
trivially data-parallel---perfect for SIMT. Diagonal gates (like the QAOA cost
layer) are just an element-wise phase multiply.
\end{analogy}

\subsection{QAOA}

The \textbf{Quantum Approximate Optimization Algorithm} is a hybrid loop for
approximately solving a QUBO/Ising problem with cost Hamiltonian
$H_C = \sum h_i Z_i + \sum J_{ij} Z_i Z_j$:
\begin{enumerate}[nosep]
  \item Start in the uniform superposition over all $2^n$ bit patterns
        (Hadamard on every qubit).
  \item Repeat $p$ times: apply $e^{-i\gamma H_C}$ (a diagonal phase: each
        amplitude $\psi_b$ is multiplied by $e^{-i\gamma\,\text{cost}(b)}$),
        then the mixer $e^{-i\beta\sum_i X_i}$ (an $R_X(2\beta)$ rotation on
        every qubit).
  \item Compute the expected cost $\langle\psi|H_C|\psi\rangle
        = \sum_b |\psi_b|^2 \text{cost}(b)$.
  \item A classical optimizer (e.g., COBYLA) adjusts the $2p$ angles
        $(\gamma,\beta)$ to lower that expectation; go to 1.
  \item Sample bit patterns from the final state; the best one is the
        answer.
\end{enumerate}
The quality metric is the \textbf{approximation ratio}: the obtained objective
divided by the true optimum (which we get by brute force).

\begin{pitfall}
Simulating QAOA on a GPU is \emph{slower} than brute-forcing the same QUBO
on the same GPU---the simulation itself touches all $2^n$ amplitudes many
times. The quantum backend in this project is not there for speed. It is
there to (a) prove the pipeline's module-detection interface can host a
quantum solver unchanged, and (b) measure how good QAOA's answers are on a
real biological graph, which matters if and when real quantum hardware is
used.
\end{pitfall}

\subsection{cuQuantum and friends}

\textbf{NVIDIA cuQuantum} is a set of CUDA libraries for simulating quantum
circuits fast on GPUs. It is a library, not a quantum computer.
\begin{itemize}[nosep]
  \item \textbf{cuStateVec:} stores $\psi$ in device memory and applies gates
        with optimized kernels (the butterflies above), computes expectation
        values, and samples measurements. This is what we use.
  \item \textbf{cuTensorNet:} represents the circuit as a network of small
        tensors and contracts them, which can handle more qubits for
        shallow or weakly-entangled circuits by never materializing $\psi$.
  \item \textbf{Frontends:} you rarely call cuQuantum directly.
        \textbf{CUDA-Q} (NVIDIA's C++/Python quantum programming model) and
        Qiskit Aer can dispatch to cuQuantum as their GPU backend. You write
        the circuit; the frontend emits cuStateVec calls.
\end{itemize}

%======================================================================
\section{Graph Neural Networks, briefly}

A \textbf{Graph Neural Network} learns a vector (embedding) for each node by
repeatedly averaging its neighbors' vectors and passing the result through a
small neural network layer. A graph convolution layer is essentially
$H' = \sigma(\hat{A} H W)$: a sparse matrix (the graph) times a dense feature
matrix times a weight matrix---SpMM followed by GEMM. Clustering the learned
embeddings gives modules. This is an optional backend; the point is that it
plugs into the same interface.

%======================================================================
\section{Where each concept lives in the pipeline}
\label{sec:map}

\begin{center}
\small
\begin{tabularx}{\linewidth}{@{}lXXl@{}}
\toprule
Module & What it does biologically & What it is computationally & Hardware \\
\midrule
0 Input & Read counts per gene per donor & Load integer matrix + metadata tables & CPU \\
1 Preprocessing & Normalize, remove batch effects, pick informative genes & Column scaling, log, linear regression, row standardization & CPU / GPU \\
2 Correlation & Which genes move together & $R=ZZ^\top$ (SYRK/GEMM) & GPU, compute-bound \\
3 TOM & Robust similarity via shared neighbors & $|R|^\beta$, row sums, $A\cdot A$, element-wise TOM & GPU, mixed \\
4 Module detection & Group genes into modules & Hierarchical clustering, or QUBO/QAOA, or GNN & CPU / GPU / cuQuantum \\
5 Biomarker output & Which modules/genes track AD & SVD per module, correlations, statistics & CPU / GPU \\
\bottomrule
\end{tabularx}
\end{center}

%======================================================================
\section{Glossary}
\label{sec:glossary}

\begin{description}[style=nextline,leftmargin=1.5em,itemsep=2pt]
  \item[Adjacency ($A$)] Edge-weight matrix of the co-expression network, $|r|^\beta$.
  \item[AMP-AD] Accelerating Medicines Partnership for AD; its consortium
        published consensus AD co-expression modules we compare against.
  \item[Braak stage / CERAD] Standard pathology scores for tangle spread and
        plaque density.
  \item[Bulk RNA-Seq] Sequencing RNA from a tissue chunk; averages over cell types.
  \item[CPM / TMM / VST] Normalizations that make samples comparable
        (counts-per-million, trimmed mean of M-values, variance-stabilizing
        transform).
  \item[DLPFC] Dorsolateral prefrontal cortex, the brain region sampled in ROSMAP.
  \item[Dynamic Tree Cut] Algorithm that cuts a dendrogram into modules at
        variable heights.
  \item[Eigengene] First principal component of a module; its summary profile.
  \item[Enrichment] Test of whether a gene list contains more members of a
        known category than chance predicts.
  \item[GEO] NCBI Gene Expression Omnibus, a public expression-data archive.
        GSE = a dataset (series), GSM = one sample, GPL = the platform.
  \item[Hub gene] Highly connected gene central to a module.
  \item[kME] Correlation of a gene with its module eigengene (module membership).
  \item[Modularity ($Q$)] Score of how well a partition places edges inside groups.
  \item[QAOA] Hybrid quantum--classical algorithm for approximate combinatorial optimization.
  \item[QUBO] Minimize $x^\top Qx$ over binary $x$.
  \item[ROSMAP] Two linked AD cohorts with brain donation; our data source.
  \item[Scale-free] Degree distribution following a power law; used to choose $\beta$.
  \item[Soft threshold ($\beta$)] Exponent turning correlations into edge weights.
  \item[State vector] The $2^n$ complex amplitudes describing $n$ qubits.
  \item[Synapse] Sage Bionetworks' data platform hosting the AD Knowledge Portal.
  \item[TOM] Topological Overlap Matrix; similarity based on shared neighbors.
  \item[WGCNA] Weighted Gene Co-expression Network Analysis, the method we accelerate.
\end{description}

\section*{Suggested next reading}
\begin{itemize}[nosep]
  \item Langfelder \& Horvath (2008), \emph{WGCNA: an R package for weighted
        correlation network analysis}, BMC Bioinformatics 9:559---the
        method paper; its tutorials walk through every step above in R.
  \item Farhi, Goldstone \& Gutmann (2014), arXiv:1411.4028---original QAOA paper.
  \item Lucas (2014), \emph{Ising formulations of many NP problems}---how to
        write problems as QUBOs.
  \item The NVIDIA cuQuantum and CUDA-Q documentation, starting with the
        cuStateVec overview.
\end{itemize}

\end{document}
